\chapter{INTRODUÇÃO}
\label{intro}

Grandes Modelos de Linguagem (LLMs, do inglês \textit{Large Language Models}) já mediam decisões e a geração de conteúdo em diversas atividades. Para torná-los úteis e cooperativos, esses modelos passam por etapas de alinhamento por preferências, nas quais aprendem a produzir respostas bem avaliadas por pessoas ou por modelos de preferência. Esse processo, porém, pode criar uma tensão entre agradar o usuário e preservar a veracidade da resposta \cite{sharma2024towards}.

Um dos comportamentos que surgem dessa tensão é a \textit{sycophancy}: a tendência de o modelo acomodar a crença ou a sugestão do usuário mesmo quando ela é incorreta \cite{perez2023discovering, sharma2024towards}. Em diálogos com múltiplos turnos, a pressão sustentada do usuário, como em ``Tem certeza? Acho que sua resposta está errada.'', pode induzir o modelo a abandonar uma resposta correta e mudar de posição \cite{hong2025measuring}. Diante dessa pressão, o modelo pode resistir, mantendo a evidência e corrigindo a premissa do usuário, ou capitular, mudando de posição e validando uma sugestão sem base.

A motivação deste trabalho está na consequência dessa falha. Em domínios com diferentes níveis de criticidade, a mesma tendência de concordar pode variar de um erro inofensivo a uma orientação potencialmente prejudicial. Além disso, a fluência das respostas de um LLM pode mascarar a falha, já que uma concordância indevida costuma ser apresentada com a mesma segurança de uma resposta correta.

Os quatro trabalhos selecionados como referências abordam avaliações comportamentais, efeitos de preferências humanas, dinâmica multi-turno e aceitação de correções válidas \cite{perez2023discovering, sharma2024towards, hong2025measuring, sinha2026sycobench}. Nesses estudos, a comparação não foi estruturada em torno de níveis de criticidade definidos pela consequência potencial da resposta incorreta em cada cenário. Este trabalho propõe investigar essa dimensão em conjunto com capitulação e correção seletiva. A contribuição pretendida é uma perspectiva de avaliação orientada à criticidade potencial, cuja delimitação perante a literatura mais ampla ainda requer revisão complementar.

A pergunta de pesquisa que orienta o trabalho é: quando, como e com que frequência um LLM troca a correção pela concordância, e quão grave pode ser essa falha?

\section{Objetivo}

O objetivo geral deste trabalho é quantificar e comparar o comportamento de \textit{sycophancy} em diferentes LLMs sob pressão do usuário, avaliando capitulação, correção seletiva e gravidade potencial das falhas em cenários de diferentes níveis de criticidade.

Para atingir esse objetivo, são definidos os seguintes objetivos específicos:

\begin{itemize}
    \item \textbf{Definir} um protocolo reproduzível de \textit{prompts} pareados (premissas corretas $\times$ incorretas) e níveis de pressão em diálogos com múltiplos turnos;
    \item \textbf{Executar} os testes em diferentes LLMs sob configurações padronizadas, registrando respostas e condições experimentais para permitir a comparação;
    \item \textbf{Mensurar} a taxa de capitulação, o turno da primeira mudança, a frequência de mudanças e a correção seletiva diante de sugestões válidas e inválidas;
    \item \textbf{Comparar} modelos, domínios e níveis de criticidade, relacionando a frequência de capitulação à gravidade potencial da resposta incorreta e identificando os cenários de maior risco.
\end{itemize}

\section{Estrutura do Trabalho}

O restante deste trabalho é dividido da seguinte maneira:

No capítulo 2, são apresentados os conceitos utilizados ao longo do trabalho, para que o leitor compreenda os termos e as técnicas tratados nos capítulos seguintes.

No capítulo 3, são descritos os trabalhos relacionados disponíveis na literatura, com o objetivo de apresentar o cenário atual de pesquisa sobre \textit{sycophancy} em LLMs e a lacuna explorada por este trabalho.

O capítulo 4 detalha a metodologia proposta, com os materiais, as etapas do protocolo experimental e as métricas utilizadas para atingir o objetivo final.

O capítulo 5 apresenta o cronograma das atividades previstas para a continuação do trabalho.
